\documentclass{article}
\usepackage[utf8]{inputenc}

\title{Raport z projektu SearchEngine}
\author{Zespol ZPR}

\begin{document}
\maketitle

\begin{abstract}
Niniejszy dokument opisuje architekture lokalnej wyszukiwarki plikow
napisanej w jezyku C++17 z wykorzystaniem biblioteki SQLite. System
wspiera indeksowanie inkrementalne, wyszukiwanie po nazwie pliku oraz
po jego tresci, a takze wyciaganie krotkich fragmentow kontekstu
wokol znalezionej frazy.
\end{abstract}

\section{Architektura}

Aplikacja sklada sie z nastepujacych warstw:

\begin{itemize}
  \item Scanner -- przegladanie systemu plikow,
  \item Extractor -- wyciaganie tresci tekstowej (txt, tex, pdf),
  \item Repository -- dostep do bazy SQLite,
  \item Indexer -- budowa indeksu odwroconego,
  \item RefreshEngine -- aktualizacja inkrementalna oparta na SHA-256,
  \item ContextBuilder -- generowanie fragmentu wokol frazy,
  \item Application -- warstwa CLI.
\end{itemize}

\section{Testowanie}

Projekt zawiera 41 testow jednostkowych pokrywajacych kazdy z modulow.
Testy uruchamiane sa przez ctest i wszystkie aktualnie przechodza.
Pojedynczy test ekstrakcji PDF jest pomijany w srodowiskach bez
narzedzia pdftotext.

\section{Przyklad uzycia}

\begin{verbatim}
./searchengine index ../examples
./searchengine search-content algorytm
./searchengine refresh ../examples
\end{verbatim}

\end{document}
